% --- 1. INFORMAZIONI GENERALI ---
\section{Informazioni Generali}
\begin{itemize}
    \item \textbf{Luogo/Piattaforma:} Server Discord\textsubscript{G}
    \item \textbf{Data:} 2026-07-14
    \item \textbf{Orario di inizio:} 14:30
    \item \textbf{Orario di fine:} 15:05
    \item \textbf{Responsabile\textsubscript{G}:} Anton Ricardo Rupi
    \item \textbf{Scriba:} Sofia De Blasi
\end{itemize}

\vspace{0.5cm}
\textbf{Partecipanti presenti:}
\begin{table}[ht]
    \centering
    \renewcommand{\arraystretch}{1.2}
    \begin{tabularx}{\textwidth}{X l}
        \toprule
        \textbf{Nome e Cognome} & \textbf{Matricola} \\
        \midrule
        Sofia De Blasi & 2111014 \\
        Roberto Mariano Doroftei & 2111031 \\
        Armando Moda Scarati & 2082864 \\
        Anton Ricardo Rupi & 2054304 \\
        Filippo Idiometri & 2035617 \\
        \bottomrule
    \end{tabularx}
\end{table}

\vspace{0.5cm}
\textbf{Partecipanti assenti:}
\begin{table}[ht]
	\centering
	\renewcommand{\arraystretch}{1.2}
	\begin{tabularx}{\textwidth}{X l}
		\toprule
		\textbf{Nome e Cognome} & \textbf{Matricola} \\
		\midrule
		Nessuno & - \\
		\bottomrule
	\end{tabularx}
\end{table}


% --- 2. ORDINE DEL GIORNO ---
\section{Ordine del Giorno}

Gli argomenti previsti per la riunione odierna sono:

\begin{enumerate}
    \item Sprint review e retrospective.
    \item Pianificazione dello sprint e assegnazione ruoli.
\end{enumerate}


% --- 3. RESOCONTO E DISCUSSIONE ---

\section{Resoconto della Riunione}

\subsection{Revisione dello sprint e retrospettiva}
Il gruppo ha effettuato la revisione dello sprint appena concluso, constatando che le attività previste sono state completate correttamente. Per quanto riguarda l’analisi dei requisiti, è necessario colmare il gap numerico e revisionare sia la coerenza complessiva sia la definizione dei requisiti.
Per quanto riguarda la retrospective, è emerso che il lavoro collaborativo su uno stesso documento rende difficile garantire la coerenza complessiva e svolgere una revisione completa ed efficace. Per questo motivo, si prevede di inserire nello sprint successivo un task dedicato, assegnato a una persona terza, con l’obiettivo di verificare e migliorare la coerenza dell’intero documento.
Tale criticità dovrà inoltre essere formalizzata come un nuovo rischio, da includere sia nello sprint sia nell’analisi dei rischi, aggiornando di conseguenza anche la relativa tabella di riepilogo.


\subsection{Pianificazione dello sprint successivo e assegnazione dei ruoli}
Il gruppo ha definito la suddivisione delle attività per lo sprint successivo, assegnando compiti specifici ai membri in base alle necessità operative emerse. Armando si occuperà dell’aggiornamento del sito e del readme, nonché della stesura dei template relative al PB. Francesco si occuperà dell’aggiornamento del PdP e del PdQ, della redazione dell’email all’azienda Zucchetti relativa al verbale esterno del 06/07/2026 e della configurazione delle automazioni CI/CD.
Ricardo curerà il verbale interno, mentre Roberto si occuperà di studiare e abbozzare design e architettura. Sofia si occuperà di rivedere l'AdR aggiornato lo sprint precedente.


% --- 4. DECISIONI PRESE ---
\section{Decisioni Prese}

\renewcommand{\arraystretch}{1.3}
\vspace{-0.3cm}
\noindent
\begin{longtable}{c p{12cm}}
    \toprule
    \textbf{Codice} & \textbf{Decisione} \\
    \midrule
    \endfirsthead

    \toprule
    \textbf{Codice} & \textbf{Decisione} \\
    \midrule
    \endhead

    \midrule
    \multicolumn{2}{r}{\textit{Continua nella pagina successiva}} \\
    \midrule
    \endfoot

    \bottomrule
    \endlastfoot

    

    \textbf{VI-22.1} & Aggiornamento del readme nella directory docs/PB, con l’aggiunta della cartella alla struttura generale e alla suddivisione per fasi di progetto. \\

    \textbf{VI-22.2} & Invio email all'azienda Zucchetti per far firmare il verbale esterno dell'incontro del 07/06 \\

\end{longtable}

% --- 5. AZIONI (TODO) ---
\section{Azioni da Svolgere (To-Do)}

\renewcommand{\arraystretch}{1.3}
\begin{longtable}{p{2.3cm} c p{4.5cm} p{3cm} l}
    \toprule
    \textbf{Codice} & \textbf{Rif. Decisione} & \textbf{Descrizione Task\textsubscript{G}} & \textbf{Assegnatario} & \textbf{Scadenza} \\
    \midrule
    \endfirsthead

    \toprule
    \textbf{Codice} & \textbf{Rif. Decisione} & \textbf{Descrizione Task} & \textbf{Assegnatario} & \textbf{Scadenza} \\
    \midrule
    \endhead

    \midrule
    \multicolumn{5}{r}{\textit{Continua nella pagina successiva}} \\
    \midrule
    \endfoot

    \bottomrule
    \endlastfoot

    \ghissue{292} & - & setup CI/CD & Francesco Marcon & 2026-07-21 \\[4pt]

    \ghissue{296} & - & Stesura verbale interno del 14/07 & Anton R. Rupi & 2026-07-21 \\[4pt]

    \ghissue{297} & VI-22.2 & Invio email per firma verbale esterno & Francesco Marcon & 2026-07-21 \\[4pt]

    \ghissue{298} & - & Aggiornamento del Piano di Progetto & Francesco Marcon & 2026-07-21 \\[4pt]

    \ghissue{299} & - & Aggiornamento del Piano di Qualifica & Francesco Marcon & 2026-07-21 \\[4pt]

    \ghissue{300} & - & Revisione dell'Analisi dei Requisiti & Sofia De Blasi & 2026-07-21 \\[4pt]   

    \ghissue{301} & - & Stesura template specifica tecnica e manuale utente & Armando M. Scarati & 2026-07-21 \\[4pt]

    \ghissue{302} & - & Studio dell'architettura e design & Roberto M. Doroftei & 2026-07-21 \\[4pt]

    \ghissue{303} & VI-22.1 & Aggiornamento del sito e readme della fase PB & Armando M. Scarati & 2026-07-21 \\[4pt]
\end{longtable}

% --- 6. PROSSIMA RIUNIONE ---
\section{Prossima Riunione}
\begin{itemize}
	\item \textbf{Data:} 2026-07-21
	\item \textbf{Ora:} 14:30
	\item \textbf{OdG preliminare\textsubscript{G}:}
	\begin{itemize}
		\item Verifica dei task assegnati per lo sprint corrente.
		\item Discussione sull'esito della comunicazione con Zucchetti.
		\item Pianificazione dello sprint successivo e assegnazione dei ruoli.
	\end{itemize}
\end{itemize}

